%% ma: L11-B09-0013
%% lop: 11
%% chuong: 3
%% bai: 9
%% dang: 11.9.4
%% loai: TL
%% muc: TH
%% dap_an: "-"
%% nguon: "de-on-tap-giua-ki-1-toan-11-nam-2026-2027-truong-thpt-ngo-quyen-da-nang.pdf (M09, Đề ôn tập giữa kì 1 Toán 11 THPT Ngô Quyền, Đà Nẵng), Mã đề 0112, Phần IV Câu 3"
%% kien_thuc: "Bài 9 (các số đặc trưng đo xu thế trung tâm)"
%% trang_thai: chinh-thuc
\begin{ex}
Khảo sát thời gian tự học (giờ/tuần) của các học sinh lớp 11B thu được bảng số liệu sau:
\begin{center}\begin{tabular}{|l|c|c|c|c|c|}\hline
Thời gian (giờ/tuần) & $[0;4)$ & $[4;8)$ & $[8;12)$ & $[12;16)$ & $[16;20]$\\\hline
Số học sinh & $4$ & $9$ & $15$ & $10$ & $2$\\\hline
\end{tabular}\end{center}
Tính thời gian tự học trung bình mỗi tuần của học sinh lớp 11B.
\loigiai{$n=40$, giá trị đại diện $2;6;10;14;18$. \[\bar x=\dfrac{2\cdot4+6\cdot9+10\cdot15+14\cdot10+18\cdot2}{40}=\dfrac{388}{40}=9{,}7\text{ (giờ/tuần).}\]}
\end{ex}
